\begin{textcolorbox}[Prompt for Rephrasing Proofs to Bound Problems]
\textbf{Task:} Transform the given inequality problem into a bound prediction problem by introducing a constant C and determining its optimal value.
\\\\
\textbf{Instructions:}

1. Analyze the original problem, focusing on its structure and potential for transformation. \\
2. Introduce a constant $C$ by either replacing an existing constant or creating a new relationship between expressions. \\
3. Determine whether to find the minimal or maximal value of $C$ that satisfies the inequality for all relevant variables. \\
4. Consider factors such as homogeneity, existing constraints, and the domain of variables (e.g., positive reals, all reals). \\
5. Ensure the rephrased problem maintains the mathematical essence and constraints of the original. \\

\textbf{Output Format:}

Provide your response in the following structure:

\texttt{<Analysis>}: Concise explanation of key features and transformation approach. \\
\texttt{<Conclusion>}: YES or NO, followed by a brief summary of the transformation. \\
\texttt{<Rephrased Problem>}: Transformed problem statement, focusing on finding the optimal $C$. \\
\texttt{<Answer>}: $C = <value>$.
\\\\
\textbf{Key Considerations:} \\
1. For double inequalities, attempt to rephrase as a single bound prediction problem when possible. \\
2. In homogeneous inequalities, focus on the ratios between variables rather than their absolute values. \\
3. Incorporate any existing constraints into the rephrased version of the problem. \\
4. Clearly specify the domain of the variables in the rephrased problem statement. \\
5. Ensure that the rephrased problem is logically equivalent to the original.
\\\\
\textbf{Example:} \\
% \texttt{\{examples\}}
Original problem:
Let $a, b, c \in \mathbb{R}^{+}$. Prove the inequality
$$
\frac{a b c}{(1+a)(a+b)(b+c)(c+16)} \leq \frac{1}{81}
$$

\vspace{2mm}
\texttt{<Analysis>}: To turn this into a bound prediction problem, we can focus on the following steps: \\
1. The left side is a rational expression that is always positive for $a, b, c \in \mathbb{R}^{+}$. \\
2. The right side is a fixed constant $\frac{1}{81}$.\\
3. We replace the constant $\frac{1}{81}$ with a variable $C$ and ask: What is the smallest $C$ such that the inequality holds for all positive $a, b, c$?\\
4. This approach allows us to determine the tightest possible upper bound for the left-hand expression.\\
5. If we find the smallest $C$ that works, we prove the original inequality and show it's the best possible.

\vspace{2mm}
\texttt{<Conclusion>}: YES, the inequality can be rephrased as a bound prediction problem. By replacing the constant $\frac{1}{81}$ with a variable $C$, we can determine the tightest upper bound for the given rational expression, effectively proving the original inequality and demonstrating its optimality.

\vspace{2mm}
\texttt{<Rephrased problem>}: \\
Coll  Determine the minimal constant $C$ such that the following inequality holds for all $a, b, c$:
$$
\frac{abc}{(1+a)(a+b)(b+c)(c+16)} \leq C.
$$
\texttt{<Answer>:} $C = \frac{1}{81}$.
\\\\
\textbf{Now, please rewrite the following problem:}

Original problem: \texttt{\{problem\}}
\end{textcolorbox}